\chapter{Discusión}\label{cap:discusion}

Las dos partes del cuerpo empírico (la ausencia de evidencia de transferencia en el
Capítulo~\ref{cap:parte1} y la atribución del retorno neto en el
Capítulo~\ref{cap:parte2}) admiten una lectura conjunta sobre dónde está, y dónde no, el
retorno en la selección transversal de acciones de gran capitalización a horizonte
semanal. La discusión se organiza en cuatro secciones. La Sección~\ref{sec:d-simples}
explica por qué los modelos simples rinden mejor en este régimen. La
Sección~\ref{sec:d-compuesto} explica por qué no es plausible que una mejora aislada
recupere el coeficiente de información. La Sección~\ref{sec:d-benchmarks} compara las
cifras del trabajo con las reportadas en la literatura. La
Sección~\ref{sec:d-limitaciones} acota el alcance de las conclusiones.

\section{Por qué los modelos simples rinden mejor}\label{sec:d-simples}

El resultado central de la Parte~I es que el IC no crece con la complejidad del modelo. Un
Lasso de cinco parámetros alcanza un IC de $+0{,}0238$, LightGBM se queda en $+0{,}0136$ y
Stockformer, con del orden de un millón de parámetros, cae a $-0{,}003$. Que el modelo más
expresivo sea el peor no es una anomalía de este panel; es, en el dominio de la selección
transversal de acciones, un fenómeno documentado por \textcite{gu2020empirical} y
confirmado en datos recientes por \textcite{rahimikia2025machine}, ya que la ganancia del
aprendizaje automático en valoración de activos procede sobre todo de no linealidades
suaves y de la regularización, no de la profundidad arquitectónica.

La explicación es un argumento clásico de sesgo y varianza, aplicado a la geometría del
problema. La unidad de información estadística no es la observación diaria, sino la
sección transversal. En cada corte temporal, el modelo tiene del orden de 400
observaciones efectivas de las que extraer una ordenación, y esas secciones están además
muy correlacionadas en el tiempo, así que el número de muestras independientes es todavía
menor. Ajustar un modelo de en torno a un millón de parámetros sobre una dimensión
efectiva tan estrecha es situarse en el régimen de alta varianza, donde el modelo tiene
capacidad de sobra para memorizar el ruido idiosincrático de cada sección transversal, y
eso es compatible con un IC fuera de muestra nulo o negativo. El Lasso está en el extremo
opuesto, ya que su penalización $L_1$ impone un sesgo fuerte que selecciona apenas cinco
coeficientes; sacrifica capacidad expresiva a cambio de una varianza muy baja, y en un
régimen con pocos datos en la dimensión que importa, ese canje sale a favor.

Este argumento es específico del régimen, no una condena universal de las redes profundas.
Stockformer se calibró sobre un mercado con ineficiencias minoristas persistentes, no
linealidades explotables y artefactos de microestructura que elevan el IC alcanzable
(Sección~\ref{sec:d-benchmarks}); en ese entorno, la capacidad adicional encuentra señal
que ajustar. El segmento de gran capitalización estadounidense es el régimen inverso, un
mercado eficiente y arbitrado donde la señal explotable es tenue y la sección transversal
estrecha. La profundidad rinde donde hay señal abundante que aprender, y penaliza donde la
señal es tenue.

\section{El efecto compuesto de los modos de fallo}\label{sec:d-compuesto}

Un lector escéptico podría objetar que el resultado de Stockformer se debe a un único
factor corregible (una pérdida mal elegida, un grafo estático, unas variables pobres o un
protocolo de validación optimista) y que la corrección adecuada recuperaría el IC. El
análisis de la Parte~I responde en parte a esa objeción, porque la escalera muestra que ni
más variables ni más capacidad llevan la señal a territorio positivo robusto. Esta sección
explica por qué, dado que los modos de fallo no actúan en paralelo, sino en serie.

La idea de fondo es una cadena de cuellos de botella multiplicativos. Cada arreglo ataca
un modo de fallo distinto. La pérdida de ordenación alinea la función objetivo con la
métrica transversal; el grafo dinámico relaja el supuesto de relaciones estáticas; las
variables enriquecidas amplían la información de entrada; el protocolo de avance temporal
quita el sesgo de una partición fija. Pero cada arreglo solo libera el cuello de botella
que aborda y deja intactos los demás. Si la señal explotable es escasa porque el mercado
es eficiente (modo a), de poco sirve alinear la pérdida con la ordenación, porque no hay
orden genuino que ordenar mejor. Si la clase de información de precio y volumen no basta a
este horizonte (modo b), enriquecer las variables dentro de esa misma clase no añade
señal, como muestra la escalera. Si parte del IC del mercado de origen era un artefacto de
microestructura inexistente en el destino (modo c), ningún cambio arquitectónico puede
recuperar una señal que nunca fue transferible. Y si la dimensión transversal efectiva es
de unos 400 nombres frente a un millón de parámetros (modo d), todo lo anterior se ve
amplificado por el sobreajuste.

El producto de varios factores, cada uno cercano a cero, es cercano a cero. Esta es la
razón estructural por la que no es plausible que una mejora aislada recupere el IC, pues
corregir un solo factor de un producto de factores próximos a cero no cambia el resultado.
Bajo este diseño, la ausencia de señal del modelo no parece deberse a un defecto puntual
de configuración, sino a que varios mecanismos de pérdida de señal actúan a la vez, y cada
uno bastaría por sí solo para acercar el IC a cero.

\section{Comparación con los resultados de la literatura}\label{sec:d-benchmarks}

Las cifras de IC de la escalera (un máximo de $+0{,}0238$ para el Lasso) son modestas
frente a las que la literatura reciente reporta para arquitecturas profundas.
\textcite{fan2024stockmixer} comunican para StockMixer un IC del orden de $0{,}041$ sobre
NASDAQ, y \textcite{li2024master} un IC del orden de $0{,}064$ para MASTER sobre CSI~300.
A primera vista, esta brecha podría leerse como prueba de que los modelos, o el marco de
evaluación, son deficientes. Hay dos consideraciones que apuntan a otra lectura, ya que esas
cifras proceden de regímenes distintos y no son directamente comparables.

La primera es de eficiencia de mercado. El IC de MASTER se obtiene sobre el CSI~300, un
mercado dominado por inversores minoristas con ineficiencias persistentes; el de
StockMixer, sobre el NASDAQ, un universo más amplio que el de las mayores
capitalizaciones del S\&P 500. La señal explotable en esos entornos es mayor que en el
segmento de gran capitalización estadounidense. Esperar que un IC obtenido en China se
reproduzca en el S\&P 500 es dar por supuesta la transferibilidad entre mercados, que es
justo el supuesto que la Parte~I somete a prueba.

La segunda afecta a la ejecutabilidad del IC reportado. El mercado chino impone límites
diarios de fluctuación de precio. Una acción que toca su límite introduce una
autocorrelación mecánica en los retornos que infla el IC alcanzable por un modelo de
precio y volumen sin que ese IC sea necesariamente operable, porque cuando la señal querría
operar al límite, el mercado está bloqueado y la operación no se ejecuta. Una parte del IC
de los resultados de mercados con límites de precio puede ser, por tanto, un artefacto no
ejecutable. El S\&P 500 no tiene esos límites, así que un IC medido en él, aunque más
bajo, es más representativo de lo que un sistema real podría obtener. La comparación de IC
entre mercados con microestructuras distintas es una comparación entre regímenes, no una
comparación de calidad de modelos en igualdad de condiciones.

\section{Limitaciones}\label{sec:d-limitaciones}

Las limitaciones siguientes no son un descargo, sino las condiciones de validez de los
hallazgos, y la mayoría se adelantaron en la Sección~\ref{sec:met-amenazas}. Aquí se
desarrollan y se ordenan.

\paragraph{Universo de datos y sesgo de supervivencia.} El universo lo forman los
componentes actuales del S\&P 500 con cobertura completa de precios, descargados de Yahoo
Finance \autocite{yahoofinance}; no se reconstruye la composición histórica del índice ni
se incluyen los valores excluidos durante el periodo (bajas de cotización, fusiones,
quiebras). Esto introduce un sesgo de supervivencia que afecta sobre todo a la pata corta
de la estrategia, porque las peores acciones de cada momento quedan infrarrepresentadas.
Una base de investigación con historial libre de ese sesgo, como CRSP
\autocite{crsp2025stock}, evitaría el problema; con los datos disponibles, los resultados
de cartera deben leerse como condicionados a este universo. Es la limitación de datos más
importante del trabajo.

\paragraph{Sin control de reproducción en el mercado de origen.} La ausencia de
transferencia se establece evaluando esta implementación del Stockformer en el S\&P 500,
pero no se ha reproducido el IC que el modelo original reporta en los índices chinos con
la misma implementación. Sin ese control, la ausencia de señal observada no puede
separarse del todo de un posible fallo de reimplementación, aunque la auditoría de fugas y
la consistencia del marco de evaluación acotan ese riesgo. Por esta razón las conclusiones
se formulan sobre esta implementación y este diseño, no sobre el modelo en abstracto.
Reproducir el IC en el mercado chino es el control más directo que queda pendiente.

\paragraph{Una sola semilla y pérdida adaptada.} El Stockformer evaluado corresponde a una
única ejecución, sin varianza por semilla, y a una pérdida adaptada (la combinación de
ordenación descrita en la Sección~\ref{sec:met-loss} en lugar del error absoluto
original). El resultado se apoya, por tanto, en una variante y una realización del modelo;
una ablación con la pérdida original y varias semillas acotaría la incertidumbre, aunque
la magnitud del resultado (IC en torno a cero) y su sobredeterminación hacen poco probable
que una semilla distinta lo revierta.

\paragraph{Datos públicos y gratuitos.} Todo el trabajo se apoya en fuentes de datos
públicas y sin coste: precios y volumen, fundamentales fechados por publicación de la SEC
\autocite{sec2025edgar} y medidas de riesgo derivadas del precio. No se dispone de los
conjuntos de datos de pago con los que operan los fondos cuantitativos (microdatos de
libro de órdenes, datos alternativos, estimaciones de analistas en tiempo real). Estos
resultados acotan lo alcanzable con datos públicos, no lo alcanzable en términos
absolutos.

\paragraph{Universo de gran capitalización y horizonte semanal.} Tanto la ausencia de
transferencia como la atribución de construcción se establecen sobre el S\&P 500 a
frecuencia de rebalanceo semanal. Las dos elecciones definen el régimen más eficiente del
espacio de posibilidades. No puede darse por hecho que las mismas conclusiones se
mantengan en capitalizaciones medianas o pequeñas, en mercados emergentes o a horizontes
intradía, donde la señal explotable y la estructura del problema son distintas.

\paragraph{Profundidad temporal y un solo país.} La evidencia cubre del orden de ocho años
y un único mercado nacional. Ocho años contienen pocos regímenes independientes (de ahí la
incertidumbre de $\pm 0{,}49$ en el Sharpe del protocolo de avance temporal), y un solo
país impide separar lo que es propio del S\&P 500 de lo que sería un fenómeno general de
los mercados eficientes.

\paragraph{Huecos en la escalera de complejidad.} La escalera salta de los árboles
potenciados (del orden de diez mil parámetros) al Stockformer (del orden de un millón) sin
un peldaño intermedio de red neuronal superficial. La omisión no es menor, porque en el
resultado de \textcite{gu2020empirical} el mejor modelo no es un árbol ni un modelo de
atención, sino una red de tres o cuatro capas. Con un único punto en la cúspide, la
evidencia muestra que este modelo no aporta señal y que el IC no crece en el tramo de los
modelos simples, pero no permite confirmar qué haría una red superficial de capacidad
intermedia. Añadir ese peldaño es el contraste pendiente más informativo de la Parte~I.

\paragraph{Costes no contabilizados y significancia económica.} La simulación contabiliza
diferencial, impacto y comisiones, pero no el coste de préstamo de títulos de la pata
corta, que con una exposición bruta de 2 no es despreciable para los nombres más difíciles
de tomar prestados; incluirlo bajaría algo el Sharpe neto, en una magnitud acotada por
tratarse de grandes capitalizaciones líquidas. Además, la significancia económica de la
estrategia final es límite, pues un Sharpe neto de $0{,}72$ con $t=1{,}82$, mejor candidato de
una búsqueda múltiple, no se distingue con confianza de cero una vez ajustado por
multiplicidad (Sección~\ref{sec:p2-modular}). El valor del trabajo no está en haber
encontrado una estrategia rentable más allá de toda duda, sino en haber cuantificado, bajo
un protocolo controlado, dónde está la palanca del retorno: en una señal-factor simple y
en la construcción de cartera, no en la arquitectura del predictor ni en más datos de la
misma clase.
